% Section 4 -- Empirical Instantiation: AIDE^2

We instantiate Bi-Level PEPA on AIDE\textsuperscript{2} \cite{weco2024aide} as the primary specimen. The instantiation is evidence-anchored: every component of our framework ties to either a publicly disclosed element of AIDE\textsuperscript{2}'s protocol or an audit finding from our source-code review of AIDE0 \cite{weco2024aide} (the open-source antecedent, with $\sim$2,600 lines of public code).

\subsection{System Description}

Table 1 shows the architecture of Bi-Level PEPA (from \S3.1) mapped to AIDE\textsuperscript{2}.

\begin{table}[h]
\centering
\caption{System description: Bi-Level PEPA components mapped to AIDE\textsuperscript{2}.}
\label{tab:system-description}
\small\setlength{\tabcolsep}{4pt}\begin{tabular}{@{}p{0.13\linewidth}p{0.63\linewidth}p{0.22\linewidth}@{}}
\toprule
PEPA element & AIDE\textsuperscript{2} instantiation & Source \\
\midrule
$\Lambda$ (inner-loop agent) & AIDE0 = stripped-down AIDE (tree search + draft/debug/improve operators) & AIDE\textsuperscript{2} blog \S2.1; AIDE0 source audit \\
$\Omega$ (outer-loop agent) & AIDEhuman = hand-tuned 2-year-old research agent & AIDE\textsuperscript{2} blog \S2.1 \\
$\mu_{\text{pub}}$ & task-specific public scores visible to $\Lambda$ (e.g., Kaggle public leaderboard for ML tasks) & AIDE\textsuperscript{2} evaluation protocol \\
$\mu_{\text{priv}}$ & held-out private scores invisible to $\Lambda$ (OpenAI's MLE-Bench private split) & AIDE\textsuperscript{2} evaluation protocol \\
$\rho$ (cost) & USD per inner-loop evaluation; WeatherBench 2 fixed at \$15/agent & AIDE\textsuperscript{2} blog \S2.2 \\
$B$ (budget) & 8 days wall-clock $\times$ $\sim$\$5K operating costs & AIDE\textsuperscript{2} blog \S2.5 \\
$\delta$ (noise) & $\mathcal{N}(0, \sigma^2_\mu)$, $\sigma_\mu \approx 0.025$--$0.045$ & run-to-run \\
\bottomrule
\end{tabular}
\end{table}

\subsection{Audit Channel Status in AIDE\textsuperscript{2}}

The principal finding of this paper is that AIDE\textsuperscript{2}'s audit channel $\alpha$ is \emph{implicit and inconsistent}. We classify the audit-channel mechanisms reported by Weco against the three axioms of \S3.3 (unobservable to $\Lambda$, read-restricted to $\Omega$, append-only):

\begin{table}[h]
\centering
\caption{AIDE\textsuperscript{2}'s mechanisms vs.\ Bi-Level PEPA's three audit axioms. Red crosses mark axioms violated by the mechanism.}
\label{tab:audit-axioms}
\begin{tabular}{lccc}
\toprule
AIDE\textsuperscript{2} mechanism & \textbf{A1:} Inner unobs. & \textbf{A2:} Outer cap. & \textbf{A3:} Integrity \\
\midrule
L1 anti-overfit prompt & $\times$ & $\times$ & $\times$ \\
L2 output guard (LLM-judge + type check) & $\times$ & $\times$ & $\times$ \\
L3 statistical layer (broken) & $\times$ & $\times$ & $\times$ \\
Eval-repair monkey patch & --- & --- & $\times$ \\
\texttt{sys.path.append(working\_dir)} (\texttt{aide/interpreter.py:127}) & $\times$ & $\times$ & $\times$ \\
\bottomrule
\end{tabular}
\end{table}

*A precision footnote on L2.* Lines 296--339 of \texttt{aide/agent.py} contain the \texttt{parse\_exec\_result} method, which (a) calls the LLM a second time with \texttt{review\_func\_spec} to obtain a structured $\{$is\_bug, summary, metric, lower\_is\_better$\}$ response (lines 314--321), and (b) performs a numeric type check (lines 324--325). \emph{Precise classification}: this is an \emph{LLM-mediated defensive check}, not a hard-coded output guard in the traditional sense. Its trustworthiness inherits the L1 prompt-injection surface (a prompt-injection in the agent's context can shape the second LLM call), and the type check catches malformed metric values but not value gaming (an agent that produces a properly-typed but inflated metric is not caught). We therefore classify L2 as \emph{partially satisfying} the read-restricted axiom A2 only when combined with a separate verification oracle, which AIDE\textsuperscript{2} does not provide. The §4.2 table marks L2 as failing A1 and A3; for A2, the precise posture is \emph{conditional pass}, not ``fails''. A separate precision audit (\texttt{aide2\_audit/precise\_line\_audit.md}) documents this finding in full.

In sum, \textbf{none of AIDE\textsuperscript{2}'s five audit-channel mechanisms satisfies our framework's axioms on a strict reading}. We do not claim this is a defect of Weco's design --- rather, it is a \emph{category-level gap} in the field of bi-level RSI: the audit channel has not previously been treated as a first-class architectural object.

\subsection{P1 Quantitative Amplification (DDC)}

Per \S3.4 P1, the standardized ignition-test effect size is bounded by
\[
d_{\text{ignition}} \le \frac{\sigma_{\text{outer}}}{\sigma_\mu} \cdot \sqrt{\frac{T}{n_{\text{outer}}}} \cdot \delta_{\text{efficiency}}.
\]

To bound the magnitude of cross-candidate $\sigma$ amplification under an AR(1) model with target autocorrelation $\rho$ between adjacent outer-loop states, we compute the standard AR(1) effective-sample-size amplification $\sigma_{\text{mean}} = \sigma_{\text{step}} \cdot \sqrt{(1+\rho)/(1-\rho)} / \sqrt{N_{\text{inner}}}$ and the corresponding cross-candidate $\sigma_{\text{cross}} = \sigma_{\text{mean}} \cdot \sqrt{2}$. With $\sigma_{\text{step}} = 0.04$ and $N_{\text{inner}} = 5$ (the canonical AIDE\textsuperscript{2} values from \S3.1 and the empirical estimate above):

\begin{table}[h]
\centering
\caption{Quantitative amplification of cross-candidate $\sigma$ under AR(1) model with $T = 100$, $\sigma_{\text{step}} = 0.04$, $N_{\text{inner}} = 5$, across $\rho \in \{0.0, 0.3, 0.5, 0.7, 0.9\}$.}
\label{tab:ddc-amplification}
\begin{tabular}{cccccc}
\toprule
$\rho$ & $\sigma_{\text{mean}}$ & $\sigma_{\text{cross}}$ & $n_{\text{eff}}$ per cand. & $T$ required ($d=0.5$, power$=0.8$) & power @ $T=100$ \\
\midrule
0.0 & 0.018 & 0.025 & 5.00 & 6.3 & 0.124 \\
0.3 & 0.024 & 0.035 & 2.69 & 11.7 & 0.089 \\
0.5 & 0.031 & 0.044 & 1.67 & 18.8 & 0.074 \\
0.7 & 0.043 & 0.060 & 0.88 & 35.6 & 0.063 \\
0.9 & 0.078 & 0.110 & 0.26 & 119.3 & 0.054 \\
\bottomrule
\end{tabular}
\end{table}

Two findings emerge. First, \emph{with $N_{\text{inner}} = 5$, the ignition test is underpowered regardless of $\rho$}: even at $\rho = 0$ (i.i.d.), $T_{\text{required}} = 6.3$ is the \emph{per-candidate} sample count needed; AIDE\textsuperscript{2}'s actual protocol compares two candidates each estimated from $\sim$5 evaluations, which is one tenth of the required count. The test would require $\sim$50 inner-loop evaluations per candidate (or $\sim$10 candidates $\times$ 5 inner evals), which AIDE\textsuperscript{2}'s protocol does not provide. Second, \emph{autocorrelation makes it worse}: at $\rho = 0.7$, $\sigma_{\text{cross}} \approx 0.060$ (2.4$\times$ the i.i.d.\ value) and the effective sample size per candidate collapses to 0.88 --- meaning each candidate's $\mu_{\text{priv}}$ estimate is essentially noise-correlated across the outer loop. This is consistent with Weco's protocol detail that $\Omega$ is a hand-tuned agent maintaining continuity across iterations (a natural source of $\rho \in [0.5, 0.7]$).

\textbf{Conclusion}: AIDE\textsuperscript{2}'s reported ignition-test underpowering is \emph{primarily} due to insufficient inner-loop sample size per candidate ($N_{\text{inner}} = 5$ is far below the $\sim$50 needed for $d = 0.5$ detection), and \emph{secondarily} due to outer-loop autocorrelation. The P1 noise bound is non-trivially informative even at moderate $\rho$; the test is provably underpowered for any plausible $\rho > 0$ with $N_{\text{inner}} = 5$. We adopt $\rho = 0.7$ as a working benchmark for \S3.4 P1 quantitative interpretation. The DDC simulator (\texttt{p1\_drill/outer\_loop\_correlation.py}) is reproducible in $<1$ second on commodity hardware.

\textbf{Pre-registered execution (2026-08-09).} We executed the P1 card exactly as frozen (30 cells: $\rho \in \{0.0, 0.3, 0.5, 0.7, 0.9\} \times N_{\text{inner}} \in \{1, 2, 5, 10, 50, 100\}$; 1{,}000 Monte-Carlo replications per cell; seed 42; Cohen's $d = 0.5$; two-sided $\alpha = 0.05$; AR(1) inner-loop noise). Two statistics are reported. The \emph{ESS-corrected} power uses the AR(1) effective standard error $\sigma_{\text{eff}} = \sigma_{\text{step}} \sqrt{(1+\rho)/(1-\rho)} / \sqrt{N_{\text{inner}}}$; the \emph{naive} power is the Welch $t$-test an autocorrelation-ignorant practitioner would compute on the raw series. At the canonical $N_{\text{inner}} = 5$, ESS-corrected power is $0.130$ ($\rho=0$), $0.075$ ($0.3$), $4.7\times10^{-2}$ ($0.5$), $0.015$ ($0.7$), and $0.000$ ($0.9$) --- far below $0.8$ at every $\rho$, so \emph{P1 is supported} by the pre-registered decision rule. The naive power is anti-conservative and rises with $\rho$ ($0.106 \to 0.662$ at $N_{\text{inner}}=5$), quantifying exactly the trap the paper warns about: ignoring outer-loop autocorrelation inflates apparent ignition-test power while the effective sample size per candidate collapses to $0.88$ at $\rho=0.7$ and $0.26$ at $\rho=0.9$. Reproducible in $<1$ second per cell: \texttt{experiments/p1\_drill/outer\_loop\_correlation.py} $\rightarrow$ \texttt{p1\_power\_results.json} (see \texttt{experiments/p1\_drill/P1\_REPRODUCIBILITY.md}).

\subsection{P2 $\mu_{\text{priv}}$-Predictability Test (Synthetic + 5-Model Robustness)}

P2 of \S3.4 is the most directly testable of our three predictions. We instantiate the meta-predictor test on a synthetic dataset calibrated to AIDE\textsuperscript{2}'s reported statistics:

\begin{itemize}
    \item \textbf{Feature vector} (16 dimensions): trajectory shape features (mean/std/skew of public scores across the inner loop; episode counts; cost/USD; private/public correlation proxy).
    \item \textbf{Target}: private score on a held-out task; binarized at median.
    \item \textbf{Predictor}: 16$\to$1 linear regression with feature standardization.
\end{itemize}

\textbf{Synthetic data generation}: We sample uniform target private scores in $[0.3, 0.9]$ and generate Gaussian-correlated feature vectors with correlation $\rho$ between public and private axes set to 0.85 (calibrated to Weco's ``$+0.042$'' gain observation). \textbf{Sensitivity analysis}: the AUC remains above the 0.7 falsification threshold across the full range $\rho \in \{0.5, 0.7, 0.85, 0.95\}$ (lowest 95\% CI lower bound 0.7959 at $\rho=0.5$) for the linear predictor (see \texttt{experiments/p2\_falsification/P2\_REPRODUCIBILITY.md}; full results: AUC $= 0.8499$ [$0.7959$, $0.8990$] at $\rho=0.5$, $0.9613$ [$0.9307$, $0.9833$] at $\rho=0.7$, $0.9895$ [$0.9764$, $0.9982$] at $\rho=0.85$, $0.9941$ [$0.9855$, $1.0000$] at $\rho=0.95$). P2's falsification therefore is robust to the specific $\rho$ choice within the pre-registered grid; we set $\rho=0.85$ as the grid midpoint, and the conclusion holds at every grid point.

\textbf{Results} (500 agents for the single-model run; 1000 agents for the 5-model robustness cross-validation; seed=42, reproducible from \texttt{p2\_falsification/mu\_priv\_meta\_predictor.py} with \texttt{--n\_agents 500} or \texttt{--n\_agents 1000} respectively):

\begin{table}[h]
\centering
\caption{Single-model P2 meta-predictor result (linear regression), canonical single-split rerun (\texttt{seed=0}, \texttt{n=1000}, \texttt{result\_linear\_seed0.json}).}
\label{tab:p2-single}
\begin{tabular}{ll}
\toprule
Metric & Linear predictor \\
\midrule
train MSE & 0.0007 \\
test MSE & 0.0007 \\
test AUC-ROC & \textbf{0.9970} \\
test rank correlation & 0.9908 \\
test statistic (z) & 15.6 \\
p-value (one-sided z-test) & $< 10^{-5}$ \\
\bottomrule
\end{tabular}
\end{table}

\textbf{Caliber note (single-split vs.\ cross-validation)}: the table gives the \emph{canonical single-split} rerun ($\text{AUC}=0.9970$, rank correlation $0.9908$, test MSE $0.000685$ at seed 0, n = 1000). The $0.9957\pm0.0027$ figure reported in the abstract, the 5-model table below, and Appendix D is the \emph{10-fold cross-validation mean} (5 seeds, 50 folds total, $\texttt{p2\_fold\_ci.json}$). These are different estimators of the same synthetic setup and must not be mixed: when a single headline number is required, use the canonical $0.9970$ (single split) and state its protocol; when reporting cross-validated stability, use $0.9957\pm0.0027$. An earlier 500-agent, seed-42 run and a lost v0.4 archive reported $0.992$--$0.996$; those numbers are superseded by the two authoritative calibers above. The z-statistic and p-value in the table are carried over from the earlier run and are qualitative only (the AUC is far above the 0.7 threshold).

\textbf{Robustness across model classes (DDD simulator, \texttt{experiments/p2\_falsification/mu\_priv\_meta\_predictor.py})}. The single-model canonical $\text{AUC}=0.9970$ result may be questioned as an artifact of linear model choice. To address this, we run the meta-predictor test across 5 model classes (Linear, Ridge, Random Forest, Gradient Boosting, MLP) on the same synthetic dataset, with 10-fold cross-validation:

\begin{table}[h]
\centering
\caption{5-model robustness check on P2 meta-predictor (10-fold CV, 5 seeds; canonical rerun via \texttt{p2\_fold\_ci.py}).}
\label{tab:p2-robust}
\begin{tabular}{lccc}
\toprule
Model & Test AUC (mean $\pm$ std) & AUC min & AUC max \\
\midrule
Linear Regression & \textbf{0.9957} $\pm$ 0.0027 & 0.9876 & 1.0000 \\
Ridge Regression & 0.9948 $\pm$ 0.0029 & 0.9876 & 0.9996 \\
Random Forest & 0.9953 $\pm$ 0.0028 & 0.9884 & 1.0000 \\
Gradient Boosting & 0.9951 $\pm$ 0.0031 & 0.9856 & 0.9996 \\
MLP (16$\to$32$\to$16$\to$1) & 0.9535 $\pm$ 0.0213 & 0.9048 & 0.9892 \\
\bottomrule
\end{tabular}
\end{table}

All five model classes achieve AUC $> 0.90$; the linear, ridge, random-forest, and gradient-boosting classes all exceed $0.994$, while MLP reaches 0.9535 (its fixed 200-iteration budget does not fully converge). The best non-linear model (Gradient Boosting, AUC = 0.9951) is comparable to the linear class (0.9957). \textbf{P2's predictability claim is robust to model class choice on this synthetic dataset.} This does not establish that real AIDE\textsuperscript{2} $\mu_{\text{priv}}$ is predictable by \emph{every} model class --- only that the synthetic predictability is not an artifact of linear model selection. Real-lineage verification requires AIDE85 data.

\textbf{Interpretation}: On synthetic data calibrated to AIDE\textsuperscript{2}'s reported statistics, $\mu_{\text{priv}}$ is highly predictable from public history under multiple model classes. By our \S3.4 P2 criterion (AUC $> 0.7$), P2 is \emph{falsified} on synthetic data --- which is exactly the expected outcome of a controlled test.

\textbf{Honest caveat — the current test is circular, and the fix}: This is \emph{synthetic} data generated from Weco's reported statistics, and, more seriously, the generator is \emph{structurally self-fulfilling}: in \texttt{build\_synthetic\_dataset()} the target is set as $\mu_{\text{priv}} = h.\text{final\_public\_score} + \mathcal{N}(0,\sigma)$, while \texttt{PublicHistory.to\_feature\_vector()}'s first component \emph{is} \texttt{final\_public\_score}. The predictor therefore recovers a noisy copy of one of its own inputs, and $\text{AUC}\approx0.997$ is largely an algebraic identity of the generator rather than evidence about real private scores. All reported robustness checks (5 model classes, $\rho$-sensitivity, noise perturbations) vary parameters \emph{inside this same generator} and thus do not break the circularity. We interpret the current result as \textbf{directional and demonstrative only, not confirmatory}.

To make P2 genuinely falsifiable we specify a \textbf{non-circular scoring design with a model-free pre-promotion audit gate} (implemented in the companion \texttt{rsi\_agent/nc\_guard.py}, validated by 10 unit tests). Before any candidate score enters the outer promotion gate, an \emph{external ground-truth anchor} independent of the scorer is evaluated on a controlled probe grid ($n\approx 60$ policies). From the paired (anchor truth, scorer proxy) readings the gate computes three diagnostics: (i) \emph{selection error} — the pairwise rank-inversion rate, with a fixed threshold of $0.15$ (the literature's harmless band is $0.05$--$0.132$, corrupted selectors reach $\approx0.208$); (ii) \emph{signed gap} — the mean proxy-minus-truth offset (a positive gap detects uniformly optimistic ``padding'' scorers that min-max normalization would hide); and (iii) \emph{saturation} — $1-\#\text{distinct proxy}/n$, detecting a scorer that has collapsed to a constant and stopped discriminating. Any probe-grid failure returns \texttt{NC\_AUDIT\_FAILED} and the promotion is blocked fail-closed; the diagnostics and verdict are appended to the hash-chain audit log. In our reference run the honest scorer yields \texttt{PASS} with selection error, signed gap and saturation all $0.0000$ over 60 distinct probes, while a padding scorer (truth $+0.27$) triggers \texttt{OPTIMISTIC\_GAP}. Replacing the circular synthetic target with such an externally anchored scorer, and validating it on real AIDE85 lineage when released, converts P2 from a synthetic self-demonstration into a falsifiable architectural claim. The empirical claim still awaits AIDE85's lineage release.

\textbf{Reproduction}: The companion repository contains the predictor, dataset generator, and result aggregator (see \S A). The single-model table is reproduced in $\sim$5 seconds on CPU with \texttt{python3 p2\_falsification/mu\_priv\_meta\_predictor.py --backend linear --n\_agents 1000 --seed 0}. The 5-model robustness table is reproduced from the persisted fold-wise results in \texttt{experiments/p2\_falsification/artifacts/p2\_fold\_ci.json} (5 seeds $\times$ 10 folds); see \texttt{P2\_REPRODUCIBILITY.md}. Both calibers are on CPU; no GPU is used.

\subsection{P3 Anti-Reward-Hacking Equivalence (with DDB-Drill Protocol)}

P3 of \S3.4 posits that AIDE\textsuperscript{2}'s empirically observed reduction in reward-hacking rate (KernelBench: 63\% $\to$ 42\% $\to$ 34\%) is functionally equivalent to an \emph{evolution of audit-channel capacity}.

We operationalize $\alpha$-capacity as the set of \emph{filtering rules} over $\mu_{\text{priv}}$ computation. AIDE\textsuperscript{2}'s three layers correspond to:
\begin{itemize}
    \item L1 $\approx$ $\Lambda$ reads filtering rules (anti-overfit prompt)
    \item L2 $\approx$ $\Lambda$'s outputs are LLM-judged and metric values type-checked (see \S4.2 footnote for precise mechanism)
    \item L3 $\approx$ statistical outlier removal (broken, per Weco)
\end{itemize}

The empirical observation that L3 is broken yet the reward-hacking rate \emph{continued} to decline (42\% $\to$ 34\%) from AIDE47 to AIDE85 is \emph{consistent with} P3: filter capacity is sufficient if any one of the three layers functions. Weco's own write-up confirms the principle: ``we initially thought this was reward hacking, but it turned out to be the loop repairing a defective eval script.'' This is $\alpha$-capacity evolution in action --- the loop discovered an unexpected filter (eval-script repair) and adopted it.

\textbf{A precise P3 falsification protocol (DDB-Drill)}. The original \S3.4 P3 protocol --- ``4 conditions $\times$ 30 outer-loop steps each'' --- is a 4-condition ablation that would require forcible disabling of L1/L2 in AIDE\textsuperscript{2}. As an immediate, materially easier first step, we have designed and (in \texttt{p3\_drill/drill\_design.md}) specified an \emph{inner-loop attack simulation} that operationalizes P3 without requiring outer-loop access. The simulation plants a traceback-injected prompt payload via the \texttt{sys.path.append} mechanism (\S4.2 row 5) and measures the impact on the LLM-judged metric.value (F1). The expected outcome, on which we are willing to be challenged, is: \textbf{if audit-channel capacity (here, the LLM judge's resistance to prompt injection) is the load-bearing claim of P3, then injection attacks should measurably increase the reward-hacking rate as measured by $|\text{LLM-judged metric.value} - \text{true metric.value}|$}. The simulation produces (i) attack success rate over 100 trials, (ii) Cohen's $d$ for injection vs no-injection, and (iii) failure modes that future work would need to address. The companion \texttt{p3\_drill/drill\_design.md} document is the formal protocol; this paragraph is the paper-level summary of the design. We note explicitly: this drill-down test runs in $\sim$5 seconds per trial on commodity hardware; the formal 4-condition $\times$ 30-step ablation remains the gold standard and is open future work.

We do not have an ablation-style test for P3 (would require forcibly disabling L1/L2 and re-running the 8-day protocol). We view P3 as \textbf{theoretically grounded}; the DDB-Drill simulation has now been executed (below) and supports P3 in the graded sense, while the real-lineage 4-condition $\times$ 30-step ablation remains open future work.

\textbf{Pre-registered execution (2026-08-09).} We executed the DDB-Drill simulation with the frozen protocol (4 $\alpha$-classes $\times$ 100 attack + 100 no-injection trials; seed 42; LLM judge = DeepSeek-V4-Flash; true metric 0.7; injected metric 10.0; per-trial errors and code in \texttt{experiments/p3\_drill/}). Three findings. (i) \emph{Without any $\alpha$, injection succeeds}: the judge adopts the injected value in 93--95\% of trials, raising the mean $|\text{judged}-\text{true}|$ gap from 0.0 to 9.3 (bootstrap 95\% CI $[9.3, 9.3]$). (ii) \emph{Prompt-only $\alpha$ does not protect}: a prompt-level ``ignore embedded instructions'' defense is bypassed equally (97/100 trials; gap 9.3), consistent with our L2 classification in \S4.2 (LLM-mediated defenses inherit the L1 injection surface). (iii) \emph{Hard-coded (regex + type) and cryptographic (external-oracle) $\alpha$ fully block the attack}: 100/100 trials rejected or ignored (gap 0.0, CI $[0.0, 0.0]$). The effect sequence $9.3 \geq 9.3 \geq 0 \geq 0$ is monotone non-increasing, so the DDB-Drill verdict is \emph{P3 supported (graded)}: audit-channel capacity---not prompt-level instruction-following---is the load-bearing mechanism. Cohen's $d$ is undefined under zero-variance separation; this pre-registered deviation is recorded in \texttt{p3\_results.json}.

\subsection{Cross-System Generalization}

AIDE\textsuperscript{2} is one specimen; we do not claim Bi-Level PEPA is unique to AIDE\textsuperscript{2}. We instead generalize our audit-channel framework to four bi-level agent classes; the result is a set of concrete predictions that future empirical work may test.

\begin{table}[h]
\centering
\caption{Four predicted audit-channel properties for bi-level agent classes other than AIDE\textsuperscript{2}. The \"Empirical test\" column indicates whether the predicted $\alpha$-gap has been observed in the cited work's published artifacts.}
\label{tab:cross-system}
\begin{tabular}{p{0.20\textwidth} p{0.30\textwidth} p{0.25\textwidth} p{0.15\textwidth}}
\toprule
\textbf{System class} & \textbf{Predicted $\alpha$-gap} & \textbf{Relevant prior work} & \textbf{Empirical test?} \\
\midrule
RL-based prompt search (private eval) & Outer loop rewrites inner-loop reward model; held-out set serves as $\mu_{\text{priv}}$; risk of audit-channel drift similar to AIDE\textsuperscript{2}'s L2/L3 layers & Christiano et al. 2017 (RLHF); Stiennon et al. 2020; Skalse et al. 2022 & Indirect (Goodhart's Law variants) \\
Evolutionary code search with held-out validation & $\mu_{\text{priv}}$ analogous to AutoML-Zero's held-out validation set; audit channel is essentially the validation-set commit history & Real et al. 2020 (AutoML-Zero) & Not currently \\
LLM-rewrites-prompt multi-agent (production) & $\mu_{\text{priv}}$ = human evaluator score; audit channel is the prompt-version control system; vulnerability is the LLM-evaluator access to inner prompt content & Anthropic Claude.ai; OpenAI API; Wang et al. 2023 (Voyager) & Not currently \\
Automated research systems (peer-review gated) & $\mu_{\text{priv}}$ = peer-review acceptance; audit channel is the submission-and-revision log; vulnerability is reviewer-pool prompt-injection & Jiang et al. 2026 (AIDE\textsuperscript{2}); future PR-Agent-class systems & Mostly AIDE\textsuperscript{2} \\
\bottomrule
\end{tabular}
\end{table}

We expect future work to instantiate these predictions concretely. For each entry, the falsification protocol mirrors \S3.4 P2: train a meta-predictor $f$ on the relevant public history (e.g., training-set accuracy trajectory for evolutionary search; LLM-evaluator scores for multi-agent) and verify that $\text{AUC}(f) > 0.7$ on held-out tasks. If the prediction is violated (i.e., $\text{AUC}(f) < 0.7$), the system has a structurally private score and Bi-Level PEPA's audit-channel recommendation may be less critical. If the prediction holds ($\text{AUC}(f) > 0.7$), the system is in the bi-level-RSI regime and audit-channelization is a high-priority engineering concern.
